\documentclass[10pt,a4paper]{amsart}
\usepackage[margin=19mm]{geometry}
\usepackage{amsmath,amssymb,mathtools}
\usepackage[scr=rsfs,cal=euler]{mathalfa}
\usepackage{xfrac}
\usepackage[unicode,hidelinks,bookmarks=false]{hyperref}
\newcommand{\ie}{i.e.\ }
\newcommand{\cf}{cf.\ }
\newcommand{\resp}{respectively\ }
\newcommand{\Subsection}{\subsection}
\providecommand{\nobreakdash}{\nobreak\hbox{-}}
\setlength{\emergencystretch}{2em}
\multlinegap0pt
\newtheorem{theorem}{Theorem}
\newtheorem{lemma}{Lemma}
\title{On Maps that Preserve the Lie Products Equal to Fixed Elements}
\author{Shiv Kumar Chaudhary, Om Prakash}
\date{Source: arXiv:2512.24912v1 (CC BY 4.0)}
\begin{document}
\maketitle

\noindent\emph{Source and licence.} On Maps that Preserve the Lie Products Equal to Fixed Elements. Version arXiv:2512.24912v1 (CC BY 4.0), distributed by arXiv under the Creative Commons Attribution 4.0 International licence, http://creativecommons.org/licenses/by/4.0/, as recorded in the arXiv licence metadata for this exact version and shown on the abstract page https://arxiv.org/abs/2512.24912v1. Full licence notice: \texttt{licenses/arxiv-2512.24912-CC-BY-4.0.txt}.\par\smallskip

\noindent\textit{Editorial scope.} Counted unit: the article's lemma lem2 (source0.tex lines 136--144). The article's complete original proof of it is delivered with it (source0.tex lines 145--149, one \texttt{proof} environment of 5 lines). No mathematics is written, repaired or re-derived: the statement and the proof are the article's own lines, in the article's own notation, and no retained line is substituted or re-flowed.  Retained uncounted as the source-local context the proof needs: lines 86--89.  Nothing was omitted from any retained span, so the package carries no omission record.  No retained line contains a citation, so the wrapper carries no bibliography and no bibliography entry.  No retained line contains a figure environment, an image-inclusion command or drawing code, so no image is delivered and the package declares no asset dependency.  Editorial formatting only, in addition: an AMS \texttt{amsart} preamble that restates the article's own declarations and macros used by the retained lines, namely the environments \texttt{theorem}, \texttt{lemma} on separate counters, together with the standard packages those lines need.  The retained environments print in this document as Theorem 1, Lemma 1 in order of appearance, whereas the article prints the same items with the numbering of its own full document.  The complete native source of the article as distributed in the arXiv e-print for this exact version is bundled unchanged as \texttt{sources/191942/source0.tex}.
\begin{theorem}
    \label{th3}
	Let $D_1 \in \mathscr{B}(\mathscr{H}_1)$ and $D_2 \in \mathscr{B}(\mathscr{H}_2)$ are fixed operators, and let $\Psi : \mathscr{B}(\mathscr{H}_1) \to \mathscr{B}(\mathscr{H}_2)$ be a linear map such that $$\Psi(A_1) \circ \Psi(A_2)=D_2 ~~\text{whenever}~~ A_1 \circ A_2=D_1$$ for every $A_1,A_2 \in \mathscr{B}(\mathscr{H}_1)$. Then $\Psi$ preserves square-zero operators.
\end{theorem}

\begin{lemma}
	\label{lem2}
	Suppose $\Psi: \mathscr{B}(\mathscr{H}_1) \to \mathscr{B}(\mathscr{H}_2)$ is a linear map satisfies the conditions of Theorem \ref{th3} if and only if it satisfies
    \begin{equation}
        \label{eq2}
        2 \Psi(X_1)^2-2\Psi(Y_1)^2=D_2~~\text{whenever}~~2X_1^2-2Y_1^2=D_1
    \end{equation}
    for all $X_1, Y_1 \in \mathscr{B}(\mathscr{H}_1)$.
\end{lemma}

\begin{proof}
	Putting $A_1=X_1+Y_1$ and $A_2=X_1-Y_1$ in the forward direction, we get
	$$\Psi(X_1+Y_1) \circ \Psi(X_1-Y_1)=D_2 ~~\text{whenever}~~ (X_1+Y_1) \circ (X_1-Y_1)=D_1,$$ and using Linearity of $\Psi$ and appropriate calculations, we get $$ 2 \Psi(X_1)^2-2\Psi(Y_1)^2=D_2 ~~\text{whenever}~~2X_1^2-2Y_1^2=D_1$$ for every $X_1, Y_1 \in \mathscr{B}(\mathscr{H}_1)$.\\
    Similarly, one can prove by substituting $X_1=\frac{A_1+A_2}{2}$ and $Y_1=\frac{A_1-A_2}{2}$ in the backward direction.
\end{proof}

\end{document}
